\newif\ifglobal

%%%%%%%%%%%%%%%%%%%%%%%
%%% Compile to view %%%
%%%%%%%%%%%%%%%%%%%%%%%

%%%%%%%%%%%%%%%%%%%%%%%%%%%%%%%%%%%%%%%%%%%%%%%%%%%
%%% Readme:                                     %%%
%%% https://www.overleaf.com/read/bwdjvfjksvjy/ %%%
%%%%%%%%%%%%%%%%%%%%%%%%%%%%%%%%%%%%%%%%%%%%%%%%%%%

% >>> M?: marks a module setting number or important notice

% >>> M4: Set {./relative-path} to external files for this module <<<
\def \modulefiles {./10-TIMG}
% To add an external file, use:
% (Replace '---' with actual filename)
% '\input{\modulefiles/tables/---.tex}' for tables
% '\includegraphics{\modulefiles/figures/---}' for figures
% '\subfile{\modulefiles/---}' for register files

% >>> M1: This module is in global project? <<<
% 'YES' - uncomment; 'NO' - comment out
\globaltrue

\ifglobal
    \documentclass[main\_\_EN.tex]{subfiles}

\else
    \documentclass[openany, 10pt]{book}

    % >>> M2: In ./00-shared/preamble-trm-module, also find
    %         settings for visibility of todo notes, watermark <<<
    \usepackage{preamble-trm-module}

    % >>> M3: Temporary labels in English,
    %         you can add your own labels there <<<
    \myexternaldocument{./00-shared/temp-labels-en}

\fi %\ifglobal


\setcounter{secnumdepth}{4}


% Define formatting of parts and chapters in text
\setenglishformatting


% >>> M5: If this module needs any updates in
% './00-shared/preamble-trm-module.sty'
% (include additional package, etc.), contact your module PM
% or create the file './00-shared/changelog.tex' make the
% necessary updates yourself and document them in this file <<<

% Make text left-aligned and allow latex to break pages naturally, comment out for full justification
\raggedright
\raggedbottom


\begin{document}
\hypertarget{timg}{}
\chapter{Timer Group (TIMG)}
\label{mod:timg}

\section{Introduction}
There are four general-purpose timers embedded in the ESP32. They are all 64-bit generic timers based on 16-bit prescalers and 64-bit auto-reload-capable up/downcounters.

The ESP32 contains two timer modules, each containing two timers. The two timers in a block are indicated by an \textcolor{red}{\it{x}} in TIMG\textcolor{red}{\it{n}}\_T\textcolor{red}{\it{x}}; the blocks themselves are indicated by an \textcolor{red}{\it{n}}.

The timers feature:
\begin{itemize}
    \item A 16-bit clock prescaler, from 2 to 65536
    \item A 64-bit time-base counter
    \item Configurable up/down time-base counter: incrementing or decrementing
    \item Halt and resume of time-base counter
    \item Auto-reload at alarm
    \item Software-controlled instant reload
    \item Level and edge interrupt generation
\end{itemize}
\section{Functional Description}

\subsection{16-bit Prescaler}
Each timer uses the APB clock (APB\_CLK, normally 80 MHz) as the basic clock. This clock is then divided down by a 16-bit precaler which generates the time-base counter clock (TB\_clk). Every cycle of TB\_clk causes the time-base counter to increment or decrement by one. The timer must be disabled (TIMG\textcolor{red}{\it{n}}\_T\textcolor{red}{\it{x}}\_EN is cleared) before changing the prescaler divisor which is configured by TIMG\textcolor{red}{\it{n}}\_T\textcolor{red}{\it{x}}\_DIVIDER register; changing it on an enabled timer can lead to unpredictable results. The prescaler can divide the APB clock by a factor from 2 to 65536. Specifically, when TIMG\textcolor{red}{\it{n}}\_T\textcolor{red}{\it{x}}\_DIVIDER is either 1 or 2, the clock divisor is 2; when TIMG\textcolor{red}{\it{n}}\_T\textcolor{red}{\it{x}}\_DIVIDER is 0, the clock divisor is 65536. Any other value will cause the clock to be divided by exactly that value.

\subsection{64-bit Time-base Counter}
The 64-bit time-base counter can be configured to count either up or down, depending on whether TIMG\textcolor{red}{\it{n}}\_T\textcolor{red}{\it{x}}\_INCREASE is set or cleared, respectively. It supports both auto-reload and software instant reload. An alarm event can be set when the counter reaches a value specified by the software.

Counting can be enabled and disabled by setting and clearing TIMG\textcolor{red}{\it{n}}\_T\textcolor{red}{\it{x}}\_EN. Clearing this bit essentially freezes the counter, causing it to neither count up nor count down; instead, it retains its value until TIMG\textcolor{red}{\it{n}}\_T\textcolor{red}{\it{x}}\_EN is set again. Reloading the counter when TIMG\textcolor{red}{\it{n}}\_T\textcolor{red}{\it{x}}\_EN is cleared will change its value, but counting will not be resumed until TIMG\textcolor{red}{\it{n}}\_T\textcolor{red}{\it{x}}\_EN is set.

Software can set a new counter value by setting registers TIMG\textcolor{red}{\it{n}}\_T\textcolor{red}{\it{x}}\_LOAD\_LO and TIMG\textcolor{red}{\it{n}}\_T\textcolor{red}{\it{x}}\_LOAD\_HI to the intended new value. The hardware will ignore these register settings until a reload; a reload will cause the contents of these registers to be copied to the counter itself. A reload event can be triggered by an alarm (auto-reload at alarm) or by software (software instant reload). To enable auto-reload at alarm, the register TIMG\textcolor{red}{\it{n}}\_T\textcolor{red}{\it{x}}\_AUTORELOAD should be set. If auto-reload at alarm is not enabled, the time-base counter will continue incrementing or decrementing after the alarm. To trigger a software instant reload, any value can be written to the register TIMG\textcolor{red}{\it{n}}\_T\textcolor{red}{\it{x}}\_LOAD\_REG; this will cause the counter value to change instantly. Software can also change the direction of the time-base counter instantly by changing the value of TIMG\textcolor{red}{\it{n}}\_T\textcolor{red}{\it{x}}\_INCREASE.

The time-base counter can also be read by software, but because the counter is 64-bit, the CPU can only get the value as two 32-bit values, the counter value needs to be latched onto TIMG\textcolor{red}{\it{n}}\_T\textcolor{red}{\it{x}}LO\_REG and TIMG\textcolor{red}{\it{n}}\_T\textcolor{red}{\it{x}}HI\_REG first. This is done by writing any value to TIMG\textcolor{red}{\it{n}}\_T\textcolor{red}{\it{x}}UPDATE\_REG; this will instantly latch the 64-bit timer value onto the two registers. Software can then read them at any point in time. This approach stops the timer value being read erroneously when a carry-over happens between reading the low and high word of the timer value.

\subsection{Alarm Generation}

The timer can trigger an alarm, which can cause a reload and/or an interrupt to occur. The alarm is triggered when the alarm registers TIMG\textcolor{red}{\it{n}}\_T\textcolor{red}{\it{x}}\_ALARMLO\_REG and TIMG\textcolor{red}{\it{n}}\_T\textcolor{red}{\it{x}}\_ALARMHI\_REG match the current timer value. In order to simplify the scenario where these registers are set 'too late' and the counter has already passed these values, the alarm also triggers when the current timer value is higher (for an up-counting timer) or lower (for a down-counting timer) than the current alarm value: if this is the case, the alarm will be triggered immediately upon loading the alarm registers. The timer alarm enable bit is automatically cleared once an alarm occurs.

\subsection{MWDT}
Each timer module also contains a Main System Watchdog Timer and its associated registers. While these registers are described here, their functional description can be found in the chapter entitled \hyperref[mod:wdt]{Watchdog Timer}.

\subsection{Interrupts}
\begin{itemize}
\item \label{int:TIMERSINTWDTINT}TIMG\textcolor{red}{\it{n}}\_INT\_WDT\_INT: Generated when a watchdog timer interrupt stage times out.
\item \label{int:TIMERSINTT1INT}TIMG\textcolor{red}{\it{n}}\_INT\_T1\_INT: An alarm event on timer 1 generates this interrupt.
\item \label{int:TIMERSINTT0INT}TIMG\textcolor{red}{\it{n}}\_INT\_T0\_INT: An alarm event on timer 0 generates this interrupt.
\end{itemize}


\hypertarget{timg-reg-summ}{}
\section{Register Summary}\label{timer_registers_summary}
\subfile{\modulefiles/timers_reg_summary_en}
{\clearpage}
\section{Registers}\label{timer_registers}
The addresses in parenthesis besides register names are the register addresses relative to the TIMG base address provided in Table \ref{table:Peripheral_Address_Mapping} \textit{Peripheral Address Mapping} in Chapter \ref{mod:sysmem} \textit{\nameref{mod:sysmem}}. The absolute register addresses are listed in Section \ref{timer_registers_summary} \textit{\nameref{timer_registers_summary}}.
\subfile{\modulefiles/timers_reg_en.tex}
\end{document}
